\documentclass[11pt,a4paper]{article}
\usepackage{fontspec}
\usepackage{kotex}
\usepackage{geometry}
\usepackage{graphicx}
\usepackage{xcolor}
\usepackage{tcolorbox}
\usepackage{booktabs}
\usepackage{array}
\usepackage{longtable}
\usepackage{colortbl}
\usepackage{fancyhdr}
\usepackage{titlesec}
\usepackage{hyperref}
\usepackage{emoji}
\usepackage{amsmath}
\usepackage{amssymb}
\usepackage{enumitem}

\geometry{left=25mm, right=25mm, top=30mm, bottom=30mm}

\setmainfont{Noto Sans CJK KR}
\setsansfont{Noto Sans CJK KR}
\setmonofont{Noto Sans CJK KR}
\setemojifont{Noto Color Emoji}

\definecolor{primary}{HTML}{1976D2}
\definecolor{darkbrown}{HTML}{3E2723}
\definecolor{mediumbrown}{HTML}{5D4037}
\definecolor{lightbrown}{HTML}{8D6E63}
\definecolor{cream}{HTML}{FFF3E0}
\definecolor{lightgray}{HTML}{F5F5F5}
\definecolor{tableheader}{HTML}{E8E8E8}

\titleformat{\section}
  {\Large\bfseries\color{primary}}
  {\thesection.}{0.5em}{}
\titleformat{\subsection}
  {\large\bfseries\color{darkbrown}}
  {\thesubsection}{0.5em}{}
\titleformat{\subsubsection}
  {\normalsize\bfseries\color{mediumbrown}}
  {\thesubsubsection}{0.5em}{}

\renewcommand{\contentsname}{목차}
\renewcommand{\figurename}{그림}
\renewcommand{\tablename}{표}

\tcbset{
  tipbox/.style={colback=cream,colframe=primary,boxrule=1.5pt,arc=3mm,
    left=5mm,right=5mm,top=3mm,bottom=3mm,fonttitle=\bfseries\color{primary}},
  formulabox/.style={colback=lightgray,colframe=lightgray,boxrule=0pt,arc=2mm,
    left=5mm,right=5mm,top=3mm,bottom=3mm},
  summarybox/.style={colback=white,colframe=primary,boxrule=2pt,arc=4mm,
    left=5mm,right=5mm,top=5mm,bottom=5mm}
}

\hypersetup{colorlinks=true,linkcolor=primary,urlcolor=primary,bookmarks=true,unicode=true}


\pagestyle{fancy}
\fancyhf{}
\fancyhead[R]{\small\color{lightbrown}마이크로비트 데이터로거 시스템 이론해설서 v1.0.0.0}
\fancyfoot[C]{\thepage}
\renewcommand{\headrulewidth}{0.4pt}

\begin{document}

\begin{titlepage}
\centering
\vspace*{3cm}
{\Huge\bfseries\color{primary}\emoji{memo} 마이크로비트 데이터로거 시스템}\\[0.5cm]
{\Huge\bfseries\color{primary} 이론해설서}
\\[1.5cm]
{\Large 시뮬레이터 버전: v1.0.3.0}
\\[2cm]
\begin{tcolorbox}[summarybox]
\centering
{\large 이 이론해설서는 마이크로비트 데이터로거 시스템이 다루는 핵심 개념과 원리를 설명합니다.}
\end{tcolorbox}
\vfill
{\large 사물인터넷과 빅데이터}\\[0.3cm]
{\large 청주교육대학교 컴퓨터교육과}
\vspace{1cm}
\end{titlepage}

\section*{1. 이 문서의 목적}

이 이론해설서는 \textbf{마이크로비트 데이터로거 시스템}이 다루는 핵심 개념과 원리를 설명합니다. 시뮬레이터를 사용하기 전에 배경 지식을 쌓거나, 사용 후 깊이 있는 이해를 위해 참조합니다.


\begin{tcolorbox}[summarybox]
\textbf{핵심 주제}: 데이터 수집과 데이터 파이프라인

\end{tcolorbox}


\section*{2. 핵심 개념}

\begin{center}
\renewcommand{\arraystretch}{1.4}
\begin{tabular}{|p{3.6cm}|p{11.4cm}|}
\hline
\rowcolor{tableheader}
\textbf{개념} & \textbf{설명} \\
\hline
\textbf{데이터 로깅} & 센서에서 읽은 값을 시각(timestamp)과 함께 저장하는 과정 \\
\hline
\textbf{RTC} & Real-Time Clock. 전원이 꺼져도 시각을 유지하는 시계 칩 \\
\hline
\textbf{CSV 파일} & Comma-Separated Values. 표 형태 데이터를 쉼표로 구분한 텍스트 파일 \\
\hline
\textbf{샘플링 주기} & 센서 데이터를 읽는 시간 간격 (예: 1초, 5초, 1분) \\
\hline
\textbf{데이터 파이프라인} & 데이터가 생성·수집·저장·분석되기까지의 전체 흐름 \\
\hline
\textbf{타임스탬프} & 데이터가 기록된 정확한 날짜와 시각 (YYYY-MM-DD HH:MM:SS) \\
\hline
\end{tabular}
\end{center}


\section*{3. 원리 상세}

\subsection*{데이터 수집의 5단계}

IoT 프로젝트의 데이터 수집은 다음 5단계로 진행됩니다:


\begin{tcolorbox}[formulabox]
\begin{verbatim}
1. 감지 → 2. 변환 → 3. 기록 → 4. 저장 → 5. 분석
\end{verbatim}
\end{tcolorbox}

\begin{center}
\renewcommand{\arraystretch}{1.4}
\begin{tabular}{|p{1.5cm}|p{7.4cm}|p{6.1cm}|}
\hline
\rowcolor{tableheader}
\textbf{단계} & \textbf{설명} & \textbf{도구/장치} \\
\hline
감지 & 물리적 현상을 전기 신호로 변환 & 센서 (SHT40, BMP280 등) \\
\hline
변환 & 전기 신호를 숫자 데이터로 변환 & ADC, I2C 프로토콜 \\
\hline
기록 & 데이터에 시각(timestamp) 추가 & RTC (DS3231) \\
\hline
저장 & 시각+데이터를 파일로 저장 & SD카드 (CSV 형식) \\
\hline
분석 & 저장된 데이터로 패턴 발견 & PC (엑셀, Python) \\
\hline
\end{tabular}
\end{center}

\subsection*{샘플링 주기의 선택}

\begin{center}
\renewcommand{\arraystretch}{1.4}
\begin{tabular}{|p{3.7cm}|p{3.0cm}|p{8.3cm}|}
\hline
\rowcolor{tableheader}
\textbf{측정 대상} & \textbf{권장 주기} & \textbf{이유} \\
\hline
온도·습도 & 1\textasciitilde{}5분 & 느리게 변하는 환경 데이터 \\
\hline
소음 수준 & 1\textasciitilde{}10초 & 비교적 빠르게 변화 \\
\hline
가속도·진동 & 10\textasciitilde{}100ms & 빠른 물리적 변화 감지 \\
\hline
기압 & 5\textasciitilde{}15분 & 매우 느리게 변화 \\
\hline
\end{tabular}
\end{center}

\textbf{주의}: 주기가 짧을수록 데이터 양이 급증합니다.

$\bullet$ 1초 간격 \texttimes{} 8시간 = 28,800건

$\bullet$ 1분 간격 \texttimes{} 8시간 = 480건



\section*{4. 실생활 응용 사례}

\textbf{사례 1}: 기상청: 전국 AWS(자동기상관측장비)가 1분 간격으로 데이터 수집


\textbf{사례 2}: 자동차 블랙박스: 가속도·GPS·영상을 초 단위로 기록


\textbf{사례 3}: 피트니스 트래커: 걸음수·심박·칼로리를 하루종일 로깅


\textbf{사례 4}: 공장 MES: 생산라인 센서 데이터를 실시간 수집·분석



\section*{5. 더 알아보기}

$\bullet$ \textbf{시계열 분석}: 시간에 따른 데이터 변화 패턴을 분석하는 통계 기법

$\bullet$ \textbf{데이터 전처리}: 결측치, 이상치 처리 등 분석 전 데이터 정리

$\bullet$ \textbf{MQTT 프로토콜}: IoT에서 널리 사용되는 경량 메시지 전송 프로토콜

$\bullet$ \textbf{클라우드 IoT}: AWS IoT, Azure IoT 등 클라우드 기반 데이터 수집·분석



\section*{6. 핵심 용어 사전}

\begin{center}
\renewcommand{\arraystretch}{1.4}
\begin{tabular}{|p{3.1cm}|p{1.7cm}|p{10.2cm}|}
\hline
\rowcolor{tableheader}
\textbf{용어} & \textbf{학술 용어} & \textbf{쉬운 설명} \\
\hline
데이터 로깅 & — & 센서에서 읽은 값을 시각(timestamp)과 함께 저장하는 과정 \\
\hline
RTC & RTC & Real-Time Clock \\
\hline
CSV 파일 & CSV 파일 & Comma-Separated Values \\
\hline
샘플링 주기 & — & 센서 데이터를 읽는 시간 간격 (예: 1초, 5초, 1분) \\
\hline
데이터 파이프라인 & — & 데이터가 생성·수집·저장·분석되기까지의 전체 흐름 \\
\hline
타임스탬프 & — & 데이터가 기록된 정확한 날짜와 시각 (YYYY-MM-DD HH:MM:SS \\
\hline
\end{tabular}
\end{center}


\section*{7. 참고자료}

\begin{center}
\renewcommand{\arraystretch}{1.4}
\begin{tabular}{|p{1.5cm}|p{6.9cm}|p{6.6cm}|}
\hline
\rowcolor{tableheader}
\textbf{\#} & \textbf{자료명} & \textbf{비고} \\
\hline
1 & 마이크로비트데이터로거\_퀵가이드 & 소프트웨어 사용 중 빠른 참조 \\
\hline
2 & 마이크로비트데이터로거\_사용설명서 & 전체 기능 상세 \\
\hline
3 & 마이크로비트데이터로거\_활용가이드 & 수업 적용 \\
\hline
4 & 마이크로비트데이터로거\_개발참조 & 기술 사양 \\
\hline
\end{tabular}
\end{center}


\vfill
\begin{center}
\small\color{lightbrown}
\textit{마이크로비트 데이터로거 시스템 이론해설서}
\end{center}
\end{document}